\chapter{Competitive Landscape / Wettbewerbslandschaft / 竞争格局}

\section{Main Competitors / Hauptkonkurrenten / 主要竞争对手}

\begin{otherlanguage}{ngerman}
Edekas Hauptwettbewerber im deutschen Markt:
\begin{itemize}
    \item \textbf{Rewe}: Zweitgrößter Player, ähnliche Positionierung
    \item \textbf{Aldi}: Discounter, starker Preisdruck
    \item \textbf{Lidl}: Discounter, internationale Expansion
    \item \textbf{Kaufland}: Hypermarket-Format
    \item \textbf{Weitere}: Real, Netto, Penny, etc.
\end{itemize}
\end{otherlanguage}

Edeka's main competitors in the German market:
\begin{itemize}
    \item \textbf{Rewe}: Second-largest player, similar positioning
    \item \textbf{Aldi}: Discounter, strong price pressure
    \item \textbf{Lidl}: Discounter, international expansion
    \item \textbf{Kaufland}: Hypermarket format
    \item \textbf{Others}: Real, Netto, Penny, etc.
\end{itemize}

艾德卡在德国市场的主要竞争对手：
\begin{itemize}
    \item \textbf{Rewe}：第二大参与者，类似定位
    \item \textbf{Aldi}：折扣店，强大的价格压力
    \item \textbf{Lidl}：折扣店，国际扩张
    \item \textbf{Kaufland}：大卖场格式
    \item \textbf{其他}：Real、Netto、Penny等
\end{itemize}

\section{Competitive Strategies / Wettbewerbsstrategien / 竞争战略}

\begin{otherlanguage}{ngerman}
Verschiedene Wettbewerbsstrategien:
\begin{itemize}
    \item \textbf{Preisführerschaft}: Aldi, Lidl
    \item \textbf{Differenzierung}: Edeka (Qualität, Nachhaltigkeit)
    \item \textbf{Nischenstrategien}: Spezialisierte Anbieter
    \item \textbf{Omnichannel}: Integration von Online und Offline
\end{itemize}
\end{otherlanguage}

% Competitive Positioning Matrix
\begin{figure}[h]
\centering
\begin{tikzpicture}[scale=1.2]
    % Axes
    \draw[->, thick] (0,0) -- (8,0) node[right] {Market Share / Marktanteil (\%)};
    \draw[->, thick] (0,0) -- (0,8) node[above, rotate=90] {Price Level / Preisniveau};
    
    % Grid
    \draw[gray!30, dashed] (0,2) -- (8,2);
    \draw[gray!30, dashed] (0,4) -- (8,4);
    \draw[gray!30, dashed] (0,6) -- (8,6);
    \draw[gray!30, dashed] (2,0) -- (2,8);
    \draw[gray!30, dashed] (4,0) -- (4,8);
    \draw[gray!30, dashed] (6,0) -- (6,8);
    
    % Labels
    \node[below] at (1,0) {\small Low};
    \node[below] at (4,0) {\small Medium};
    \node[below] at (7,0) {\small High};
    \node[left] at (0,1) {\small Low};
    \node[left] at (0,4) {\small Medium};
    \node[left] at (0,7) {\small High};
    
    % Competitors
    \node[circle, fill=blue!70, draw=blue!90, thick, minimum size=1.2cm] (edeka) at (6.5,6.5) {};
    \node[above] at (6.5,7.2) {\textbf{Edeka}};
    \node[below] at (6.5,5.8) {\small 26\%};
    
    \node[circle, fill=red!70, draw=red!90, thick, minimum size=1cm] (rewe) at (4.5,6) {};
    \node[above] at (4.5,6.7) {Rewe};
    \node[below] at (4.5,5.3) {\small 15\%};
    
    \node[circle, fill=green!70, draw=green!90, thick, minimum size=0.9cm] (aldi) at (3.5,2.5) {};
    \node[above] at (3.5,3.2) {Aldi};
    \node[below] at (3.5,1.8) {\small 12\%};
    
    \node[circle, fill=yellow!70, draw=yellow!90, thick, minimum size=0.8cm] (lidl) at (2.5,2) {};
    \node[above] at (2.5,2.7) {Lidl};
    \node[below] at (2.5,1.3) {\small 10\%};
    
    \node[circle, fill=orange!70, draw=orange!90, thick, minimum size=0.7cm] (kaufland) at (2,4.5) {};
    \node[above] at (2,5.2) {Kaufland};
    \node[below] at (2,3.8) {\small 8\%};
    
    % Quadrants
    \node[above right, gray!50] at (0,8) {Premium};
    \node[above left, gray!50] at (8,8) {Differentiation};
    \node[below right, gray!50] at (0,0) {Budget};
    \node[below left, gray!50] at (8,0) {Mass Market};
    
    % Title
    \node[above, font=\bfseries] at (4,8.5) {Competitive Positioning Matrix / Wettbewerbspositionierungs-Matrix / 竞争定位矩阵};
\end{tikzpicture}
\caption{Competitive Positioning / Wettbewerbspositionierung / 竞争定位}
\end{figure}

Different competitive strategies:
\begin{itemize}
    \item \textbf{Price leadership}: Aldi, Lidl
    \item \textbf{Differentiation}: Edeka (quality, sustainability)
    \item \textbf{Niche strategies}: Specialized providers
    \item \textbf{Omnichannel}: Integration of online and offline
\end{itemize}

不同的竞争战略：
\begin{itemize}
    \item \textbf{价格领导}：Aldi、Lidl
    \item \textbf{差异化}：艾德卡（质量、可持续性）
    \item \textbf{利基战略}：专业提供商
    \item \textbf{全渠道}：线上线下整合
\end{itemize}

\section{Market Dynamics / Marktdynamik / 市场动态}

\begin{otherlanguage}{ngerman}
\begin{itemize}
    \item Konsolidierungstrends
    \item Markteintritte und -austritte
    \item Preiskriege
    \item Innovationstrends
    \item Regulatorische Auswirkungen
\end{itemize}
\end{otherlanguage}

\begin{itemize}
    \item Consolidation trends
    \item Market entries and exits
    \item Price wars
    \item Innovation trends
    \item Regulatory impacts
\end{itemize}

\begin{itemize}
    \item 整合趋势
    \item 市场进入和退出
    \item 价格战
    \item 创新趋势
    \item 监管影响
\end{itemize}

\section{Competitive Advantages / Wettbewerbsvorteile / 竞争优势}

\begin{otherlanguage}{ngerman}
Edekas Wettbewerbsvorteile:
\begin{itemize}
    \item Größte Marktposition
    \item Starke Eigenmarken
    \item Regionale Präsenz
    \item Qualitätsimage
    \item Nachhaltigkeitspositionierung
    \item Genossenschaftsmodell (Flexibilität)
\end{itemize}
\end{otherlanguage}

Edeka's competitive advantages:
\begin{itemize}
    \item Largest market position
    \item Strong private labels
    \item Regional presence
    \item Quality image
    \item Sustainability positioning
    \item Cooperative model (flexibility)
\end{itemize}

艾德卡的竞争优势：
\begin{itemize}
    \item 最大的市场地位
    \item 强大的自有品牌
    \item 区域存在
    \item 质量形象
    \item 可持续性定位
    \item 合作社模式（灵活性）
\end{itemize}
